\chapter{Building a chat client}
\label{ch:chatclient}

\begin{aside}
A walkthrough: build a chat client TUI in \maya{} from
scratch. Three panels, async networking, message history.
\end{aside}

\section{Sketch}

Three vertical panels: channel list (left), message view
(centre), input (bottom). A status bar at the top.

\section{State}

\begin{lstlisting}
struct ChatState {
    Signal<std::vector<Channel>> channels;
    Signal<int> active_channel;
    Signal<std::vector<Message>> messages;
    Signal<std::string> input_text;
    Signal<ConnState> conn;
};
\end{lstlisting}

\section{Layout}

\begin{lstlisting}
auto root = column(
    status_bar() | height(1),
    row(
        channel_list() | flex(1),
        column(
            message_view() | flex(1),
            input_box()    | height(3)
        ) | flex(4)
    ) | flex(1)
);
\end{lstlisting}

\section{Networking}

A background task connects to the server, parses incoming
messages, posts to the message signal.

\begin{lstlisting}
spawn([&]() -> Task<void> {
    auto sock = co_await connect("chat.example");
    while (sock.alive()) {
        auto msg = co_await sock.read_message();
        ctx.post_to_main([msg] {
            messages.push_back(msg);
        });
    }
});
\end{lstlisting}

\section{Input}

The bottom input is a text input. Enter sends the message:

\begin{lstlisting}
on_key(KeyEnter, [&] {
    auto text = ctx.get(input_text);
    if (!text.empty()) {
        send(text);
        ctx.set(input_text, "");
    }
});
\end{lstlisting}

\section{Scrolling}

The message view auto-scrolls to bottom on new messages
unless the user has scrolled up.

\section{Notifications}

Highlighted messages (your name mentioned) trigger an OSC 9
notification.

\section{Persistence}

Recent messages cached to disk; reloaded on startup.

\section{Recap}

\begin{recap}
\begin{itemize}[leftmargin=*]
  \item Three panels + status bar.
  \item Async networking via spawn + post\_to\_main.
  \item Auto-scroll with user-override.
  \item Notifications on mention.
  \item Cache to disk.
\end{itemize}
\end{recap}

\section*{Workshop}

\begin{enumerate}[leftmargin=*]
  \item Build the layout shell.
  \item Wire a mock network backend.
  \item Add notifications and persistence.
\end{enumerate}
